% SPDX-License-Identifier: CC-BY-SA-4.0
% Author: Matthieu Perrin
% Part: <Nom de la partie>
% Section: <Nom de la section>
% Sub-section: <Nom de la sous-section>  % (facultatif, laisser vide si non utilisé)
% Frame: <Titre de la slide>

\begingroup

\begin{frame}{Inclusion des classes de complexité}

  \onBlock<1>[top=-3mm]{Quelques résultats inconnus}{
    \begin{minipage}{.5\textwidth}
      \begin{itemize}
      \item $\textsc{p} \stackrel{?}{=} \textsc{np}$
      \item $\textsc{np} \stackrel{?}{=} \textsc{co-np}$
      \end{itemize}
    \end{minipage}%
    \begin{minipage}{.5\textwidth}
      \begin{itemize}
      \item $\textsc{np} \stackrel{?}{=} \textsc{pspace}$
      \item $\textsc{p} \stackrel{?}{=} \textsc{np} \cap \textsc{co-np}$
      \end{itemize}
    \end{minipage}
  }

  \onBlock<1>[bottom=5mm]{Théorème de Savitch}{
    \begin{itemize}
      \item Pour toute fonction $f(n) \ge \log(n)$,
      \quad
      \alert{$\textsc{nspace}\left( f(n) \right) \subseteq \textsc{dspace}\left( f(n)^2 \right)$}
    \end{itemize}
  }
  
  \on[y=-3mm]{
    
    \begin{tikzpicture}[x=4.6mm, y=4.6mm]
      \uncover<2>{
        \draw[alert, fill=alert!20]                                     \theshape{7};
        \draw[black, fill=black!20]                       \theshape{6};
        \begin{scope}
          \fill[rotate= 7, white]                        \theshape{5};
          \fill[rotate=-7, white]                        \theshape{5};
          \fill[rotate= 7, structure!40, opacity=.5]     \theshape{5};
          \fill[rotate=-7, example!40, opacity=.5]       \theshape{5};
          \draw[rotate= 7, structure]                    \theshape{5};
          \draw[rotate=-7, example]                      \theshape{5};
        \end{scope}
        \draw[alert, fill=alert!20]                      \theshape{4};
      }
      \draw[black, fill=black!20]                        \theshape{3};
      \begin{scope}
        \fill[rotate= 17, white]                         \theshape{2};
        \fill[rotate=-17, white]                         \theshape{2};
        \fill[rotate= 17, structure!40, opacity=.5]      \theshape{2};
        \fill[rotate=-17, example!40, opacity=.5]        \theshape{2};
        \draw[rotate= 17,  structure]                    \theshape{2};
        \draw[rotate=-17, example]                       \theshape{2};
      \end{scope}
      \draw[alert, fill=alert!20]                        \theshape{1};

      \footnotesize
      \node[alert]                              at ( 0.5, 0) {$\textsc{p}$};
      \node[structure]                          at ( 2.5, 2) {$\textsc{np}$};
      \node[example]                            at ( 2.5,-2) {$\textsc{co-np}$};
      \node[structure!50!example, align=center] at ( 3.0, 0) {$\textsc{np}$ \\ $\cap$ \\ $\textsc{co-np}$};
      \uncover<2>{
        \node[alert]                              at ( 8.0, 0) {$\textsc{exptime}$};
        \node[structure]                          at ( 6.5, 5) {$\textsc{nexptime}$};
        \node[example]                            at ( 6.5,-5) {$\textsc{co-nexptime}$};
        \node[alert]                              at (15.0, 0) {$\textsc{r}$};
      }
      \scriptsize
      \node[black, align=center]                at (5.5,0)  {$\textsc{pspace}$ \\ $=$ \\ $\textsc{npspace}$ \\ $=$ \\ $\textsc{co-}$ \\ $\textsc{npspace}$};
      \tiny
      \uncover<2>{
        \node[structure!50!example, align=center] at (10.1,0) {$\textsc{nexptime}$ \\ $\cap$ \\ $\textsc{co}$ \\ $\textsc{-nexptime}$};
        \node[black, align=center]                at (12.5,0) {$\textsc{expspace}$ \\ $=$ \\ $\textsc{nexpspace}$ \\ $=$ \\ $\textsc{co-}$ \\ $\textsc{nexpspace}$};
      }
      
      \footnotesize
      \draw[black!60, densely dotted]      (-7.7,0) -- (16.5,0);
      \node[black!60, rotate=90, left ] at (-7.5,0) {symétrie par};
      \node[black!60, rotate=90, right] at (-7.5,0) {complémentaire};

    \end{tikzpicture}
  }    

  \only<1>{
    \footnoteref{W. J. Savitch. \textit{Relationships between nondeterministic and deterministic tape complexities.} JCSS (1970)}
  }

\end{frame}


\endgroup
\endinput
